
\clearpage
\VerdierPhysicalPageAnchor{0123}%
\begingroup
\small
\setlength{\emergencystretch}{1.5em}%
\setlength{\abovedisplayskip}{4pt plus 2pt minus 2pt}%
\setlength{\belowdisplayskip}{4pt plus 2pt minus 2pt}%

\VerdierSemanticAnchor{ch2-1-2-14}
\noindent\textbf{注 1.2.14.}\quad
在 1.2 中研究的性质从未用到公理 TRIV。因此，这些性质对于任何装备了一族
满足公理 TRI、TRII 与 TRIII 的三角的严格加法 $\mathbb Z$-范畴都成立。

\VerdierSemanticAnchor{ch2-1-3}
\medskip
\noindent\textbf{1.3. 三角范畴的例子。}

\VerdierSemanticAnchor{ch2-1-3-1}
\medskip
\noindent\textbf{1.3.1.}\quad
设 $C$ 为一个加法范畴，$K(C)$ 为 $C$ 的复形同伦范畴（第一章，2.5.7）。
范畴 $K(C)$ 是一个严格 $\mathbb Z$-范畴，并装备了一族三角，即特异三角
（第一章，3.3.1）；这族三角满足性质 TRI、TRII 与 TRIII（第一章，3.3）。

\VerdierSemanticAnchor{ch2-1-3-2}
\medskip
\noindent\textbf{命题 1.3.2.}\quad
\textit{范畴 $K(C)$ 连同这族特异三角构成一个三角范畴。}

\VerdierSemanticAnchor{ch2-1-3-2-triv-reduction}
只需证明 $K(C)$ 的这族特异三角具有性质 TRIV。为此，设

\VerdierSemanticAnchor{ch2-1-3-2-commutative-triangle}
\[
  \begin{tikzcd}[column sep=4.2em,row sep=2.6em]
    & X_2 \arrow[dr,"u_1"] & \\
    X_1 \arrow[ur,"u_3"] \arrow[rr,"u_2"'] && X_3
  \end{tikzcd}
\]
\noindent 为一个交换图，并设

\VerdierSemanticAnchor{ch2-1-3-2-input-triangles}
\[
  \begin{aligned}
    X_1&\xrightarrow{u_3}X_2\xrightarrow{v_3}Z_3
      \xrightarrow{w_3}X_1[1],\\
    X_2&\xrightarrow{u_1}X_3\xrightarrow{v_1}Z_1
      \xrightarrow{w_1}X_2[1],\\
    X_1&\xrightarrow{u_2}X_3\xrightarrow{v_2}Z_2
      \xrightarrow{w_2}X_1[1]
  \end{aligned}
\]
\noindent 为三个特异三角。由（第一章，3.3.5），存在三个特异三角

\VerdierSemanticAnchor{ch2-1-3-2-alternative-triangles}
\[
  \resizebox{0.98\linewidth}{!}{$
    (X_1,X_2,Z'_3,u_3,v'_3,w'_3),\qquad
    (X_2,X_3,Z'_1,u_1,v'_1,w'_1),\qquad
    (X_1,X_3,Z'_2,u_2,v'_2,w'_2)
  $}
\]
\noindent 以及两个态射

\VerdierSemanticAnchor{ch2-1-3-2-morphisms-boundary}
\[
  m'_1\colon Z'_3\longrightarrow Z'_2,\qquad
  m'_3\colon Z'_2\longrightarrow Z'_1,
\]
\endgroup


